% AI-written, AI cross-reviewed; no human sign-off; not a claim that RH is proved
% Manuscript v3, heading p. 148; original paper2/pending_proof.tex, SHA-256 f080c15d3cb18fd56ac15e23b3bc092ea515925739b4aee4e5561c8df4058a40.
\subsection{Candidate proof of depth zero in the base frame (P134)}
\label{sec:pending-p134}
\noindent\textsf{P134. The Round 21 AI reviewer found the statement and argument correct and referred the candidate for separate cross-review.}
The issue is to identify the lattice stabilized in the \emph{chosen}
base frame. An arbitrary rational frame need not have the property.

\paragraph{The local test.}
Work in the chart's coefficient field $E$, and let $Q=s(P)$ be the
transport of one marked prime. Write $\widehat O_Q$ for the valuation ring
in $E_Q$ and $V_Q=\widehat O_Q^n$, where $n$ is the stage rank.
In the source, an adelic scalar $z$ is integral at $Q$ if some
$s\in\mathcal O_E\setminus Q$ makes $sz$ integral at all finite places.
This is equivalent to $z_Q\in\widehat O_Q$. Indeed, such an $s$ is a
unit at $Q$. Conversely, an adele has denominators at only finitely many
other primes $w$. Choose $s\equiv1\pmod Q$ and
$s\equiv0\pmod{w^{e_w}}$ at those primes, taking $e_w$ large enough to
clear the corresponding denominators. The Chinese remainder theorem
gives one such $s$. Here the tensor adelic model gives the same integral
lattice: at a place $v$ of the base field $F$, completing the finite
$\mathcal O_F$-module $\mathcal O_E$ gives
$\prod_{w\mid v}\widehat O_w$, by factorization of $v\mathcal O_E$ and
the Chinese remainder theorem on its powers. A fixed finite family of
$\mathcal O_F$-module generators suffices at all $v$, so the integral
tensors generate precisely the product of these rings.
Consequently depth zero for $k$ tests the entries of $k_Q-I$ for
integrality. At depth $d$ it tests the entries of
$\mathrm N(Q)^{-d}(k_Q-I)$.

\paragraph{Identification of the base frame.}
Let $L$ be the chart level's base lattice, $J$ its Hermitian matrix and
$a$ the scalar in the tied data. Of the two hypotheses in \nolinkurl{UnderTie},
the base-frame calculation uses the first, \nolinkurl{AdPins}. Its
\nolinkurl{crd_orth} identity gives $J=H^*H$ for the invertible coordinate
matrix $H=A.\mathrm{crd}_i$, so $\det J$ is a unit. Write
$C=L+aL^{\sharp}$, where the dual is taken with respect to $J$.
This is a full lattice, also when $a=0$. Indeed $L$ is full and finite,
and the dual is obtained by testing a finite spanning family against the
nonsingular pairing. In a local basis these tests define a lattice by
an invertible matrix. Denote the completion of $C$ by $C_Q$.
For a unitary adele $x$, adelic transport sends the completed lattice at
each place to $x_wL_w$. To see that these local conditions recover the
transported global lattice, choose a common nonzero integer $h$ with
$hL_w\subset x_wL_w\subset h^{-1}L_w$ for all $w$. Such an $h$ exists
because $x$ and its inverse are integral away from finitely many places,
as is a basis change for $L$. The finite quotient $h^{-1}L/hL$ is the
direct sum of its primary parts. Specifying the submodule in each part
and taking its inverse image gives exactly the global rational vectors
satisfying the local conditions, with the prescribed completions.

Duality commutes with this transport. The defining test for the
Hermitian dual is integrality of all pairings with lattice vectors.
After completion one may test on local generators, since the global
ones span there. At a split place the paired component is at the
conjugate prime. For the full tuple of components the identity
$\langle xu,xv\rangle_J=\langle u,v\rangle_J$ carries the test in both
directions. Thus the completion of the companion of the transported
lattice at $Q$ is $x_Q C_Q$.

Put $b=A.\mathrm{fb}_{i,\beta,P}$ and let
$R$ denote the decorated point's rational frame transported to $E$.
Under the field identification a denominator outside $P$ remains outside
$Q$, so localization, and then completion, carries the decorated frame
test to $Q$. In the generator-prime case this test is part of the
point's lifted decoration. At the other marked primes it is the frame
condition of the point itself. In both cases $R$ takes the standard
lattice onto the transported decorated lattice. The lattice clause
of the represented point identifies this last lattice with the image
under $H$ of the companion of the lattice transported by $x$. Hence
\[
                  R V_Q=H x_Q C_Q .
\]
The frame clause is independent of that lattice clause. It says that
$K=(H x_Q b)^{-1}R$ is congruent to $I$ at depth $d=(t.\mathrm{cst}\,\mu).D(P)$. The index condition
\nolinkurl{(t.idx_ok μ).cst_ok.D_pos P} gives $d\ge1$.
Since the rational integer $\mathrm N(Q)$ belongs to $Q$, we have
$K\equiv I$ modulo the maximal ideal of $\widehat O_Q$.
Its determinant is a unit and the adjugate formula makes $K^{-1}$
integral. Therefore $K V_Q=V_Q$, which gives
$R V_Q=H x_Q b V_Q$. Cancelling $H x_Q$ from the two equalities yields
\[
                              b V_Q=C_Q .
\]

\paragraph{Stabilization and inverse.}
For an element $\kappa$ of the level subgroup, each local component
lies in the stabilizer of the completed base lattice, by the local
level inclusion. The same holds for $\kappa^{-1}$ because the level is a
group. The unitary-pairing test just used for $x$ shows that $\kappa$
also preserves the dual lattice, taking account of both primes in a
split pair. It therefore preserves $L+aL^\sharp$ in both directions.
At $Q$, conjugating by $b$ now gives
\[
 (b^{-1}\kappa_Q b)V_Q=V_Q,
 \qquad (b^{-1}\kappa_Q^{-1} b)V_Q=V_Q .
\]
Applying these matrices to the coordinate vectors gives integral
columns. Subtracting the corresponding columns of $I$ preserves
integrality. The local test proves the two depth-zero congruences of
\nodeid{P134}. In rank zero the entrywise assertions are vacuous.

\smallskip
\noindent\textit{Lean comparison.}
The point's lattice and frame fields are in
\nolinkurl{UnderTie.lean:24--27}. The two cases for its decorated frame
are in \nolinkurl{IntegralGLSuccEmbeddingMemParahoricTimesPrincipa.lean:275--280}.
The base-frame calculation is in
\nolinkurl{ChartBaseFrameBadPrimeEqCompanionLattice.lean:55--83}.
The stabilizer and the entrywise denominator argument are in
\nolinkurl{GlobLatCompanOfLev.lean}, under
\nolinkurl{globLat_compan_of_lev}, \nolinkurl{int_conj} and
\nolinkurl{lev_stab_cong}. The parent statement links to the full
source location. The source pins for this candidate are in
\path{out/R21_pending.json}. The review and the revised source comparisons for both candidates are recorded in \path{out/R22_candidate_review.json}.
